\section{Finite parameter consequences}\label{finite-parameter-consequences}

In this section a triple system is \emph{connected} when its Levi graph is
connected.  We work first without isolated vertices.  For an integer $r\ge1$,
put
\[
q(r)=\left\lceil2\sqrt r\right\rceil.
\]

\begin{proposition}[Edge-deletion form of the bridge condition]
\label{proposition-edge-deletion-bridge-condition}
Let $F$ be a finite triple system and let $e=\{x,y,z\}\in E(F)$.  The
hyperedge-node $e$ has no incident bridge in $I(F)$ if and only if $x,y,z$
lie in one connected component after the hyperedge $e$ is deleted and all
points are retained.  Consequently the bridge condition in Theorem~A is
equivalent to the following hypergraph-native condition:
\[
\text{for every }e\in E(F^\circ),\text{ deleting }e\text{ separates its
three points into at least two components.}
\]
\end{proposition}
\begin{proof}
Fix $x\in e$.  The Levi incidence $xe$ is not a bridge precisely when there
is an alternative path from $e$ to $x$ after that incidence is deleted.  A
simple such path leaves $e$ through $y$ or $z$ and thereafter lies entirely
in the Levi graph of $F-e$.  Thus $xe$ is non-bridging exactly when $x$ is
connected in $F-e$ to at least one of the other two points of $e$.

If all three incidences at $e$ are non-bridging, no one of $x,y,z$ can form a
singleton component among those three points.  A partition of three points
with this property has one block, so all three are connected in $F-e$.  The
converse follows by reversing the same paths.
\end{proof}

\begin{lemma}[Bipartite shadow]
\label{lemma-bipartite-shadow}
Let $F$ be an obligatory finite triple system with no isolated vertices and
at least one hyperedge.  There is a finite bipartite graph $J$, without
isolated vertices, such that
\[
|E(J)|=|E(F)|,\qquad |V(J)|=|V(F)|-|E(F)|,
\]
and $J$ has the same number of connected components as $F$.  Conversely, for
every finite bipartite graph $J$ without isolated vertices, the expansion
$J^+$ is obligatory and has these parameter relations.
\end{lemma}
\begin{proof}
By Theorem~A and Proposition~\ref{proposition-5.2-bridge-block-decomposition},
each connected component $H$ of $F$ is obtained by one-point amalgamations
from connected expansion pieces
\[
J_1^+,\ldots,J_k^+,
\]
where each $J_i$ is a connected bipartite graph.  Disconnected graph pieces
may simply be split into their connected components.  Write
$m_i=|E(J_i)|$ and $s_i=|V(J_i)|$.  Since $H$ is connected, its $k$ pieces
are joined by exactly $k-1$ one-point identifications.  Hence
\[
|E(H)|=\sum_{i=1}^k m_i,
\qquad
|V(H)|=\sum_{i=1}^k(m_i+s_i)-(k-1).
\tag{10.1}
\]

Take one-point sums of the ordinary graphs $J_1,\ldots,J_k$ along any tree on
the $k$ indices.  Before each identification, interchange the two
bipartition classes of the new factor if necessary.  The resulting graph
$J_H$ is connected and bipartite, and (10.1) gives
\[
|E(J_H)|=|E(H)|,
\qquad
|V(J_H)|=|V(H)|-|E(H)|.
\]
Taking the disjoint union of these shadows over the components of $F$ proves
the forward assertion.  The converse is the expansion-generator case of
Theorem~A, with the displayed counts immediate from the definition of $J^+$.
\end{proof}

\begin{theorem}[Exact order--size--component spectrum]
\label{theorem-order-size-component-spectrum}
Let $m\ge1$, $1\le c\le m$, and $n$ be integers.  There exists an obligatory
triple system $F$ with no isolated vertices, exactly $m$ hyperedges, exactly
$n$ vertices, and exactly $c$ connected components if and only if
\[
\boxed{
 m+2(c-1)+\left\lceil2\sqrt{m-c+1}\right\rceil
 \le n\le 2m+c.
}
\tag{10.2}
\]
Every integer in this interval occurs.
\end{theorem}
\begin{proof}
We first record the corresponding elementary graph fact.  A connected simple
bipartite graph with $r\ge1$ edges and $s$ vertices exists if and only if
\[
q(r)\le s\le r+1.
\tag{10.3}
\]
Indeed, connectedness gives $r\ge s-1$.  If the bipartition has sizes $a,b$,
then
\[
r\le ab\le\left\lfloor\frac{s^2}{4}\right\rfloor,
\]
which is equivalent to $s\ge q(r)$.  Conversely,
$K_{\lfloor s/2\rfloor,\lceil s/2\rceil}$ contains a spanning tree, and
starting with that tree and adding arbitrary unused edges realises every edge
count between $s-1$ and $\lfloor s^2/4\rfloor$.

Apply Lemma~\ref{lemma-bipartite-shadow}.  Let the $c$ connected components of
the shadow have positive edge counts $m_1,\ldots,m_c$, where
$\sum_i m_i=m$.  If their orders are $s_i$, then (10.3) gives
\[
\sum_{i=1}^c q(m_i)\le \sum_{i=1}^c s_i\le m+c.
\tag{10.4}
\]
For positive integers $a,b$,
\[
q(a)+q(b)\ge q(a+b-1)+2.
\tag{10.5}
\]
To see this, note that
\[
\sqrt a+\sqrt b\ge\sqrt{a+b-1}+1,
\]
because, after squaring, the assertion reduces to
$(\sqrt a-1)(\sqrt b-1)\ge0$; taking ceilings gives (10.5).  Repeatedly
applying (10.5) to the left side of (10.4) yields
\[
\sum_{i=1}^c q(m_i)
\ge q(m-c+1)+2(c-1).
\]
The shadow has $n-m$ vertices, so these inequalities give (10.2).

For sharpness put $M=m-c+1$.  Take $c-1$ components equal to $K_2$.  By
(10.3), a final connected bipartite component with $M$ edges can have every
order from $q(M)$ through $M+1$.  Their disjoint union therefore realises
every shadow order from $q(M)+2(c-1)$ through $m+c$.  Expanding it adds
exactly $m$ private vertices and proves attainability of every integer in
(10.2).
\end{proof}

\begin{corollary}[Connected and unrestricted spectra]
\label{corollary-connected-order-size-spectrum}
For $m\ge1$:
\begin{enumerate}
\item a connected obligatory triple system without isolated vertices has
$m$ hyperedges and $n$ vertices if and only if
\[
m+\left\lceil2\sqrt m\right\rceil\le n\le2m+1;
\]
\item without prescribing the number of components, a reduced obligatory
triple system has $m$ hyperedges and $n$ vertices if and only if
\[
m+\left\lceil2\sqrt m\right\rceil\le n\le3m;
\]
\item if isolated vertices are permitted, an obligatory triple system with
$m$ hyperedges and $n$ vertices exists if and only if
\[
n\ge m+\left\lceil2\sqrt m\right\rceil.
\]
\end{enumerate}
\end{corollary}
\begin{proof}
The first assertion is Theorem~\ref{theorem-order-size-component-spectrum}
with $c=1$.  For the second, the connected interval covers through $2m+1$;
for $n>2m+1$, take $c=n-2m$ and use the upper endpoint $n=2m+c$ of (10.2).
The largest possible value is obtained from $m$ disjoint triples.  The final
assertion follows by adjoining arbitrary isolated vertices.
\end{proof}

\begin{corollary}[Exact Levi cycle-rank spectrum]
\label{corollary-levi-cycle-rank-spectrum}
Let $F$ be reduced and obligatory, with $m$ hyperedges, $n$ vertices, and $c$
connected components.  The cyclomatic number of its Levi graph is
\[
\beta(I(F))=|E(I(F))|-|V(I(F))|+c=2m-n+c.
\tag{10.6}
\]
For fixed $m,c$, every integer in the interval
\[
0\le\beta(I(F))\le
m-c+2-\left\lceil2\sqrt{m-c+1}\right\rceil
\tag{10.7}
\]
occurs, and no other value occurs.  In particular,
\[
|V(F)|=2|E(F)|+c
\quad\Longleftrightarrow\quad
I(F)\text{ is a forest}.
\]
\end{corollary}
\begin{proof}
The Levi graph has $n+m$ vertices, $3m$ incidence edges, and $c$ components,
which gives (10.6).  Substitution of the two endpoints in (10.2) gives
(10.7), and every intermediate value occurs because every intermediate order
does.  A finite graph has cyclomatic number zero exactly when it is a forest.
\end{proof}

\begin{corollary}[Rigidity at balanced lower endpoints]
\label{corollary-balanced-endpoint-rigidity}
Let $F$ be connected, reduced, and obligatory.
\begin{enumerate}
\item If $|E(F)|=t^2$ and $|V(F)|=t^2+2t$, then
$F\cong K_{t,t}^+$.
\item If $|E(F)|=t(t+1)$ and $|V(F)|=t(t+1)+2t+1$, then
$F\cong K_{t,t+1}^+$.
\end{enumerate}
Here $t\ge1$, and the complete bipartite graphs are unique up to interchanging
the two classes.
\end{corollary}
\begin{proof}
In either case the shadow from Lemma~\ref{lemma-bipartite-shadow} attains
$|E(J)|=\lfloor|V(J)|^2/4\rfloor$.  Equality forces every possible cross edge
to be present and the two bipartition classes to be balanced, so the shadow is
$K_{t,t}$ or $K_{t,t+1}$, respectively.

For $t\ge2$ these complete bipartite graphs have no cut vertex.  On the other
hand, the shadow built from two or more connected expansion pieces is a
nontrivial one-point sum and has a cut vertex.  Hence the bridge-block
reconstruction has only one piece, and $F$ is the expansion of the displayed
complete bipartite graph.  The cases $t=1$ are immediate.
\end{proof}

